\section{Conditional two-family entropy estimate}

\begin{theorem}[Labelled conditional information bound]
    \label{thm:peps_labelled_conditional_entropy}
    \lean{TNLean.PEPS.AreaLaw.regionalEntropy_conditional_le_residual_add_sum}
    \leanok
    \uses{def:area_law_ground_entropy, def:area_law_two_family_partition}
    Let $\Omega$ be a normalized pure state on a finite domain, and let $G,B,C$
    be pairwise-disjoint regions. Suppose an ordered two-family partition
    (Definition~\ref{def:area_law_two_family_partition}) of $G$ has remainder $D$,
    pieces $X_i$, family labels $f_i\in\{0,1\}$, and preceding family unions
    $P_i=\bigcup_{j<i,\,f_j=f_i}X_j$.
    If $I_\Omega(X_i:C\cup B\cup P_i)\le\varepsilon_i$ when $f_i=0$ and
    $I_\Omega(X_i:(G\cup B)^c\cup P_i)\le\varepsilon_i$ when $f_i=1$, then
    (\cite[Lemma~3.2]{OpenAI2026PolynomialPEPS})
    \begin{align}
        S_\Omega(G\cup B)+S_\Omega(C\cup B)-S_\Omega(B)-S_\Omega(G\cup C\cup B)
        \le 2S_\Omega(D)+\sum_i\varepsilon_i.
        \label{eq:peps_labelled_conditional_entropy}
    \end{align}
\end{theorem}
\begin{proof}\leanok
    Filtering the common order by $f_i\in\{0,1\}$ identifies each family's
    preceding union with $P_i$. Since $C\cap(G\cup B)=\varnothing$,
    $C\cup(G\cup B\cup C)^c=(G\cup B)^c$, giving the $f_i=1$ hypothesis
    the complementary-exterior form; the generic pure-state
    conditional two-family bound and disjoint index sums then yield
    (\ref{eq:peps_labelled_conditional_entropy}), as finite-domain regional entropies
    (Definition~\ref{def:area_law_ground_entropy}) agree with generic reduced-state entropies.
\end{proof}

% Source: OpenAI2026PolynomialPEPS, Lemma 3.2 (lem:cell-cmi), 02-information.tex,
% lines 197--222.
\begin{corollary}[Counting form of the conditional information bound]
    \label{cor:peps_cell_information_counting_form}
    \lean{TNLean.PEPS.AreaLaw.regionalEntropy_conditional_le_card_residual_add_pieceCount}
    \leanok
    \uses{thm:peps_labelled_conditional_entropy}
    In the setting of Theorem~\ref{thm:peps_labelled_conditional_entropy}, if
    the local dimension is $q$, the partition has $N$ pieces, and
    $\varepsilon_i=C_2$ for all $i$, then
    \begin{align}
        S_\Omega(G\cup B)+S_\Omega(C\cup B)-S_\Omega(B)-S_\Omega(G\cup C\cup B)
        \le 2\lvert D\rvert\log q+NC_2.
        \label{eq:peps_cell_information_counting_form}
    \end{align}
\end{corollary}
\begin{proof}\leanok
    Substitute $\sum_{i=1}^N\varepsilon_i=NC_2$ and $S_\Omega(D)\le\lvert D\rvert\log q$
    (from $\operatorname{rank}\Omega_D\le q^{\lvert D\rvert}$) into (\ref{eq:peps_labelled_conditional_entropy}).
\end{proof}

% Source: OpenAI2026PolynomialPEPS, Lemma 3.2 (lem:cell-cmi), 02-information.tex,
% lines 310--311.
\begin{corollary}[Information of a small cell]
    \label{cor:peps_small_cell_information}
    \lean{TNLean.PEPS.AreaLaw.regionalEntropy_conditional_le_two_mul_card_mul_log}
    \leanok
    \uses{cor:peps_cell_information_counting_form}
    Let $\Omega$ be a normalized pure state on a finite domain with local
    dimension $q$, and let $G,B,C$ be pairwise-disjoint regions. Then
    \begin{align}
        S_\Omega(G\cup B)+S_\Omega(C\cup B)-S_\Omega(B)-S_\Omega(G\cup C\cup B)
        \le 2\lvert G\rvert\log q.
        \notag
    \end{align}
\end{corollary}
\begin{proof}\leanok
    Apply Corollary~\ref{cor:peps_cell_information_counting_form} with $N=0$ pieces and $D=G$.
\end{proof}

% Source: OpenAI2026PolynomialPEPS, Lemma 3.2 (lem:cell-cmi), 02-information.tex,
% lines 224--227, 240--241, 276--280 and 302--311.
\begin{theorem}[Generation count of the cell recursion]
    \label{thm:peps_cell_generation_count}
    \lean{TNLean.PEPS.Approximation.cellChildSide,
        TNLean.PEPS.Approximation.exists_cellChildSide_le_rpow,
        TNLean.PEPS.Approximation.exists_cell_generation_bound,
        TNLean.PEPS.Approximation.exists_cell_generation_count_bound}
    \leanok
    \uses{thm:area_law_power_recursion_depth}
    In the single-cell recursion, an active square of side $u$ has children
    of side $u'=8u^{\gamma}$. For $\gamma<\gamma'$, there exists $u_*>1$ with
    $8u^{\gamma}\le u^{\gamma'}$ for all $u\ge u_*$. Given $0<\gamma'<1$ and any
    such $u_*>1$, there exists $C>0$ such that for every initial side $s>0$ and
    every $m$ whose generations $0,\dots,m-1$ have sides $\ge u_*$,
    \begin{align}
        m & \le C+C\log\log(s+e).
        \label{eq:peps_cell_generation_depth}
    \end{align}
    Moreover, for every $b\ge0$ there exists $C>0$ such that, under the same condition with $s\ge1$,
    \begin{align}
        \sum_{j=0}^{m}b^{j} & \le C(1+\log s)^{C}.
        \label{eq:peps_cell_generation_sum}
    \end{align}
\end{theorem}
\begin{proof}\leanok
    For $u\ge u_*\ge8^{1/(\gamma'-\gamma)}$, $8u^{\gamma}\le u^{\gamma'}$, so
    $0<u'\le u^{\gamma'}$ and Theorem~\ref{thm:area_law_power_recursion_depth} with $\rho=\gamma'$
    yields (\ref{eq:peps_cell_generation_depth}). For $s\ge1$,
    $\log(s+e)\le\log(e^{2}s)=2+\log s\le2(1+\log s)$, yielding (\ref{eq:peps_cell_generation_sum}).
\end{proof}

% Source: OpenAI2026PolynomialPEPS, Lemma 3.2 (lem:cell-cmi), 02-information.tex,
% lines 152--158 and 224--300.
\begin{lemma}[One active square of the cell recursion]
    \label{lem:peps_cell_active_square}
    \lean{TNLean.PEPS.Approximation.CellPartition.IsChamberSquare,
        TNLean.PEPS.Approximation.CellPartition.holes,
        TNLean.PEPS.Approximation.CellPartition.tile,
        TNLean.PEPS.Approximation.CellPartition.IsChamberSquare.colour_eq_of_ne,
        TNLean.PEPS.Approximation.CellPartition.IsChamberSquare.lt_dist_of_mem_tile,
        TNLean.PEPS.Approximation.CellPartition.IsChamberSquare.isRectUnion_tileSample,
        TNLean.PEPS.Approximation.CellPartition.IsChamberSquare.isChamberSquare_child}
    \leanok
    Label every point of $\mathbb R^2$ as allowed or forbidden of colour $c\in\{0,1\}$,
    with chambers and squares following one half-open boundary convention.
    Let $Q$ be a square of side $u$ with at most $5$ vertical and $5$ horizontal lines
    (including $\partial Q$), with distinct parallel lines separated by $\ge u/8$,
    such that labels are constant on their chambers in the $u$-neighbourhood of $Q$.
    For $0<t\le u/1000$, reserve the radius-$4t$ square about every line intersection
    within distance $10t$ of $Q$, and let $Q^\circ$ be the unreserved part of $Q$. Then:
    \begin{enumerate}
        \item every $p\in Q^\circ$ sees within distance $t$ at most one label different from its own;
        \item setting tile $T_1=\{p\in Q^\circ:\text{$p$ allowed},\,\operatorname{dist}(p,\text{colour }0)\le t\}$
              and $T_0=\{p\in Q^\circ:\text{$p$ allowed}\}\setminus T_1$, every
              $p\in T_i$ satisfies $\operatorname{dist}(p,\text{colour }i)>t$;
        \item for $i\in\{0,1\}$, $T_i\cap\mathbb Z^2$ is a union of at most $66^2$
              closed lattice rectangles;
        \item each reserved square, with lines at offsets $0,\pm t,\pm4t$ from its centre,
              satisfies the same hypothesis with side $8t$ once each $T_i$ is relabelled with colour $i$.
    \end{enumerate}
\end{lemma}
\begin{proof}\leanok
    A label change within distance $t$ of $p\in Q^\circ$ requires a line within distance $t$ of $p$.
    Two distinct parallel lines within $t$ of $p$ would be $\le 2t<u/8$ apart, while orthogonal lines
    would intersect within $t$ of $p$, placing $p$ inside a radius-$4t$ reserved square. Thus at most one
    line lies within $t$ of $p$, proving (1) and hence (2). Clamping a point near $p$ into the
    $t$-ball of $q$ preserves chambers whenever $p,q$ share a chamber of the lines shifted by $\pm t$;
    hence the distance test in (2) is constant on chambers of the old lines, their $\pm t$-shifts,
    and reserved-square sides, yielding (3). Within $12t$ of a reserved square, only the two lines
    through its centre and no other reserved square appear, so the updated labels are constant on
    chambers of the lines at offsets $0,\pm t,\pm4t$, proving (4).
\end{proof}

% Source: OpenAI2026PolynomialPEPS, Lemma 3.2 (lem:cell-cmi), 02-information.tex,
% lines 181--188 and 227--300.
\begin{lemma}[Generations of the cell recursion]
    \label{lem:peps_cell_generation_invariant}
    \lean{TNLean.PEPS.Approximation.CellPartition.Generation.next,
        TNLean.PEPS.Approximation.CellPartition.Generation.IsValid,
        TNLean.PEPS.Approximation.CellPartition.Generation.IsValid.next,
        TNLean.PEPS.Approximation.CellPartition.initialGeneration_isValid,
        TNLean.PEPS.Approximation.CellPartition.pairwise_pieceList}
    \leanok
    \uses{lem:peps_cell_active_square}
    A generation consists of a common side $u$, finitely many active squares of
    side $u$ at mutual distance $>10u$ satisfying the hypothesis of
    Lemma~\ref{lem:peps_cell_active_square}, and labels such that every allowed point
    lies in an active square. Initially, a single allowed grid cell of side $s$ with
    all exterior cells forbidden of colour $0$ or $1$ forms a generation.
    If $1000u^{\gamma}\le u$, colouring each tile $T_i$ (at $t=u^{\gamma}$) with $i\in\{0,1\}$
    and taking the reserved squares as active squares of side $u'=8u^{\gamma}$ produces
    a valid next generation. Ordered by generation and then by active square, the tiles
    are pairwise disjoint, and every $p\in T_i$ has distance $>u^{\gamma}$ from every
    earlier tile of family $i$.
\end{lemma}
\begin{proof}\leanok
    Children of one parent have centres $\ge u/8>80u^{\gamma}$ apart, while children of distinct
    parents lie within $14u^{\gamma}$ of parents $>10u$ apart, so child separation exceeds $10u'=80u^{\gamma}$.
    Lemma~\ref{lem:peps_cell_active_square}(4) preserves the chamber invariant since only a parent's own tiles
    change labels near it, and remaining allowed points lie in reserved squares. Once coloured $i$,
    a tile $T_i$ is avoided by later family-$i$ tiles by Lemma~\ref{lem:peps_cell_active_square}(2),
    while same-generation tiles in distinct active squares are $>10u>u^{\gamma}$ apart.
\end{proof}

% Source: OpenAI2026PolynomialPEPS, Lemma 3.2 (lem:cell-cmi), 02-information.tex,
% lines 181--196 and 224--311.
\begin{theorem}[Two-family partition of a cell]
    \label{thm:peps_cell_two_family_partition}
    \lean{TNLean.PEPS.Approximation.CellPartition.exists_cell_twoFamily_partition}
    \leanok
    \uses{lem:peps_cell_generation_invariant, thm:peps_cell_generation_count,
        def:area_law_two_family_partition}
    Fix $\gamma<1$ and $u_0$. There exist $u_*\ge\max(u_0,1)$ and constants $C,K>0$
    such that for any axis-parallel square grid of side $s\ge1$ with one chosen cell
    and all exterior cells coloured $0$ or $1$, the lattice sample of the chosen cell
    admits an ordered two-family partition (Definition~\ref{def:area_law_two_family_partition})
    whose pieces $X_j$ (of family $f_j\in\{0,1\}$) each have a scale $u_j\in[u_*,s]$ such that:
    \begin{enumerate}
        \item $X_j$ is a union of at most $66^2$ closed lattice rectangles of $\ell_\infty$-diameter $\le u_j$;
        \item every $x\in X_j$ has $\ell_\infty$-distance $>u_j^{\gamma}$ from every lattice point of any
              exterior cell of colour $f_j$ and from every $X_i$ with $i<j$ and $f_i=f_j$;
        \item the total number of pieces plus remainder sites is at most $C(1+\log s)^{K}$.
    \end{enumerate}
\end{theorem}
\begin{proof}\leanok
    Choose $0<\gamma<\gamma'<1$ and $u_*\ge\max(u_0,1)$ so that $1000u^{\gamma}\le u$ and
    $8u^{\gamma}\le u^{\gamma'}$ for all $u\ge u_*$. Iterating Lemma~\ref{lem:peps_cell_generation_invariant}
    until $u<u_*$ terminates in $m$ generations bounded by (\ref{eq:peps_cell_generation_depth}) of
    Theorem~\ref{thm:peps_cell_generation_count}. Collecting all tiles gives the pieces $X_j$ and
    the remaining allowed sites give $D$. Since each active square spawns $\le25$ children, generation $g$
    has $\le25^g$ active squares, $2$ tiles per square, and $\le(u_*+1)^2$ lattice sites per terminal square;
    (\ref{eq:peps_cell_generation_sum}) with $b=25$ bounds the total by $C(1+\log s)^K$.
\end{proof}
